\chapter*{Preface}
\addcontentsline{toc}{chapter}{Preface}
\markboth{Preface}{Preface}

The idea of having a virtual lab came to me during the discussions with my wife,
who is a development scientist in the dairy industry. She was explaining different
methods of experimental design that come up during their discussions with her
colleagues, and how to analyze the data, especially when the number of
experiments is limited. You can imagine that a lot of these discussions become speculative until
the actual experiments are conducted and the data analyzed. If the method is a
failure, a lot of resources are wasted. I was thinking about some sort of virtual
laboratory: one that is rooted in the actual mechanisms governing the system under
investigation. If we can have this mechanistic model, perhaps with some stochastic
elements in it, we can run the model with the input parameters that come from an
actual DoE method. A mechanistic stochastic model can even give us the variability
we can expect in the real world, and the uncertainty we can mimic by adding noise
to the virtual ``measurements''. The measurements are another story! We rarely have any
direct measurement of the desired outcome of an experiment: when we need
concentration, we measure a peak. When we need to know the stage of a fermentation
reaction, we measure pH or turbidity. So it is important to have mechanistic models
that can give us these measured values. During my wife's PhD, she had a special
course with me on the stochastic models for bacterial growth, so it was relatively
easy for me to create something based on the formulations I had collected. For the
downstream biological separation processes, I had a similar model for the
mass-transfer-controlled oxygen scavenging process using MOFs. The project was with
my friend Jonas and I wrote a small Julia package for it with a formulation that I
was very proud of at the time. I fed all these codes and documents (that are in
\TeX{} format) to Claude Code and asked it to draft a Python package for me. I was
not very happy with the models because of some simplifications in the mass transfer
model into porous particles. I first used a general mass transfer model and tried
to solve it with the method of lines, but the solver failed (on second thought, I
should have used a solver suited for stiff ODEs). I therefore switched to our PyFVTool
finite volume solver. I gave it a sample coupled problem I had solved for
multiphase flow in porous media. Claude Code, however, really impressed me when it
managed to use PyFVTool to couple two domains---the chromatographic column and the
particles---and solve the nonlinear system of PDEs. I gradually extended the
package based on the DoE methods I had heard before and those that my wife
explained to me. Eventually, I prompted Claude to convert each module to a chapter
in a monograph and add two foundation chapters to briefly explain the concepts.
I meant this part for light reading during a train ride, or as a fast track to
familiarize potential users with these fascinating concepts. Here we are now! I
hope you like the product and send me your comments for improvements. I will come
back to it if my time permits.

\bigskip
\noindent\hfill\emph{Ali A. Eftekhari}
